\documentclass[a4paper,11pt]{article}

% ============================================================
%  Le Collé — Algèbre linéaire (compléments de deuxième année)
%  © 2026 Ahmed Amine Ouqlani
% ============================================================

% ---------- Polices ----------
\usepackage[utf8]{inputenc}
\usepackage[T1]{fontenc}
\usepackage[sc]{mathpazo}          % texte et mathématiques : Palatino
\usepackage[scaled=0.88]{helvet}   % titres et repères : linéale
\linespread{1.07}
\usepackage{amsmath,amssymb,amsthm,mathrsfs,stmaryrd}
\usepackage{microtype}

% ---------- Mise en page ----------
\usepackage[a4paper,top=2.3cm,bottom=2.3cm,left=2.4cm,right=2.4cm,
            headheight=14pt,headsep=0.8cm,footskip=1.1cm]{geometry}
\setlength{\parindent}{0pt}
\setlength{\parskip}{0.45em plus 0.15em minus 0.1em}
\raggedbottom
\emergencystretch=1.5em
\clubpenalty=10000 \widowpenalty=10000 \displaywidowpenalty=10000 \brokenpenalty=10000

% ---------- Couleurs : un bleu de nuit, un ocre, des gris ----------
\usepackage{xcolor}
\definecolor{Ink}{HTML}{1C2533}
\definecolor{Bleu}{HTML}{1F3A5F}
\definecolor{BleuMoyen}{HTML}{4A6A8F}
\definecolor{BleuClair}{HTML}{9DB2C9}
\definecolor{Fond}{HTML}{F2F4F7}
\definecolor{Ocre}{HTML}{B07A26}
\definecolor{Brique}{HTML}{A23B2A}
\definecolor{Gris}{HTML}{5E6672}
\definecolor{Filet}{HTML}{C5CCD6}

% ---------- Packages ----------
\usepackage{enumitem}
\setlist{itemsep=0.25em,topsep=0.35em,parsep=0pt}
\setlist[itemize]{label=\textcolor{BleuMoyen}{\small\textbullet},leftmargin=1.4em}
\setlist[enumerate]{leftmargin=1.8em,label=\textsf{\small\color{Bleu}\arabic*.}}
\usepackage{booktabs,array,multicol}
\usepackage{needspace,float}
\usepackage[font=small,labelfont={sf,color=Bleu},labelsep=period,skip=6pt]{caption}
\usepackage{tikz}
\usetikzlibrary{arrows.meta,calc,decorations.markings}
\usepackage{pgfplots}
\pgfplotsset{compat=1.17}
\usepackage{tkz-tab}
\usepackage[most]{tcolorbox}
\usepackage{titlesec,titletoc}
\usepackage{fancyhdr,lastpage,etoolbox}
\usepackage[colorlinks=true,linkcolor=Bleu,urlcolor=Bleu,citecolor=Bleu,
            pdftitle={Le Collé — DL 3 : bases et familles cycliques},pdfauthor={Ahmed Amine Ouqlani}]{hyperref}

\numberwithin{equation}{section}
\numberwithin{figure}{section}
\numberwithin{table}{section}

% ---------- Titres ----------
% numéro du chapitre dans le parcours PT (affiché), distinct de la numérotation interne
\newcommand{\grandnum}{\thesection}\newcommand{\ongletcouleur}{Bleu}
% Chapitres (sections) : grand numéro, titre en linéale, filet ocre
\titleformat{\section}[display]
  {\normalfont\sffamily\color{Ink}}
  {{\small\color{Gris}\textls[200]{\MakeUppercase{Chapitre}}}\\[-0.2em]%
   {\fontsize{64}{64}\selectfont\color{BleuClair}\grandnum}}
  {0.4em}
  {\fontsize{25}{30}\selectfont\bfseries}
  [\vspace{0.5em}{\color{Ocre}\rule{3.2cm}{1.4pt}}]
\titlespacing*{\section}{0pt}{0pt}{1.6em}
\titleformat{name=\section,numberless}[display]
  {\normalfont\sffamily\color{Ink}}{}{0pt}{\fontsize{25}{30}\selectfont\bfseries}
  [\vspace{0.5em}{\color{Ocre}\rule{3.2cm}{1.4pt}}]
\titlespacing*{name=\section,numberless}{0pt}{0pt}{1.6em}

\titleformat{\subsection}
  {\needspace{5\baselineskip}\normalfont\sffamily\Large\bfseries\color{Ink}}
  {\color{Bleu}\thesubsection}{0.8em}{}
\titlespacing*{\subsection}{0pt}{1.8em}{0.5em}
\titleformat{\subsubsection}
  {\needspace{5\baselineskip}\normalfont\sffamily\large\bfseries\color{Ink}}
  {\color{Bleu}\thesubsubsection}{0.7em}{}
\titlespacing*{\subsubsection}{0pt}{1.6em}{0.4em}
\titleformat{\paragraph}[runin]{\normalfont\sffamily\bfseries\color{Ink}}{}{0pt}{}[.]
\titlespacing*{\paragraph}{0pt}{0.9em}{0.6em}
\setcounter{secnumdepth}{3}
\setcounter{tocdepth}{2}

% ---------- Table des matières ----------
\titlecontents{section}[0pt]{\addvspace{0.7em}\sffamily\bfseries\color{Ink}}
  {}{}{\hfill\contentspage}[\addvspace{0.2em}]
\titlecontents{subsection}[2.2em]{\small\linespread{1}\selectfont}
  {\makebox[2.6em][l]{\color{Gris}\thecontentslabel}}{}{\titlerule*[0.7pc]{.}\contentspage}

% ---------- En-têtes et pieds de page ----------
\pagestyle{fancy}
\fancyhf{}
\renewcommand{\sectionmark}[1]{\markboth{#1}{}}
\fancyhead[L]{\footnotesize\sffamily\color{Gris}\leftmark}
\fancyhead[R]{\footnotesize\sffamily\scshape\color{Gris}Le Collé}
\fancyfoot[C]{\small\sffamily\color{Ink}\thepage}
\fancyfoot[L]{\scriptsize\sffamily\hyperlink{toc}{\color{Filet}sommaire}}
\renewcommand{\headrulewidth}{0.4pt}
\renewcommand{\headrule}{{\color{Filet}\hrule height\headrulewidth width\headwidth}}
\renewcommand{\footrulewidth}{0pt}
\fancypagestyle{plain}{\fancyhf{}\fancyfoot[C]{\small\sffamily\color{Ink}\thepage}\renewcommand{\headrulewidth}{0pt}}

% ---------- Énoncés ----------
\newtheoremstyle{enonce}{0pt}{0pt}{}{}{\sffamily\bfseries\color{Bleu}}{.}{0.5em}
  {\thmname{#1}\thmnumber{~#2}\thmnote{\ \normalfont\sffamily\color{Ink}(#3)}}
\newtheoremstyle{exemple}{0.9em}{0.9em}{}{}{\sffamily\bfseries\color{BleuMoyen}}{.}{0.5em}
  {\thmname{#1}\thmnumber{~#2}\thmnote{\ \normalfont\sffamily\color{Ink}(#3)}}
\newtheoremstyle{remarque}{0.7em}{0.7em}{}{}{\sffamily\itshape}{.}{0.5em}{\thmname{#1}\thmnote{\ (#3)}}
\theoremstyle{enonce}
\newtheorem{definition}{Définition}[section]
\newtheorem{theoreme}[definition]{Théorème}
\newtheorem{proposition}[definition]{Proposition}
\newtheorem{corollaire}[definition]{Corollaire}
\theoremstyle{exemple}
\newtheorem{exemple}[definition]{Exemple}
\theoremstyle{remarque}
\newtheorem*{remarque}{Remarque}

\tcbset{
  defi/.style={enhanced,breakable,colback=white,colframe=white,boxrule=0pt,arc=0pt,
     borderline west={1.2pt}{0pt}{Bleu},left=11pt,right=4pt,top=3pt,bottom=3pt,
     before skip=1.1em,after skip=1.1em},
  thm/.style={enhanced,breakable,colback=Fond,colframe=Fond,boxrule=0pt,arc=0pt,
     borderline west={2.4pt}{0pt}{Bleu},left=11pt,right=8pt,top=6pt,bottom=6pt,
     before skip=1.2em,after skip=1.2em}}
\tcolorboxenvironment{definition}{defi}
\tcolorboxenvironment{theoreme}{thm}
\tcolorboxenvironment{proposition}{thm}
\tcolorboxenvironment{corollaire}{thm}

\renewcommand{\proofname}{Démonstration}
\renewcommand{\qedsymbol}{\textcolor{BleuClair}{$\blacksquare$}}
\makeatletter
\renewenvironment{proof}[1][\proofname]{\par\pushQED{\qed}\normalfont
  \topsep0.5em\relax\trivlist\item[\hskip\labelsep\sffamily\itshape\color{Gris}#1.]\ignorespaces}
  {\popQED\endtrivlist\@endpefalse}
\makeatother

% ---------- Méthode : filets haut et bas, titre en bandeau ----------
\makeatletter\newcommand{\methlabelfix}{\protected@edef\@currentlabel{\thetcbcounter}}\makeatother
\newtcolorbox[auto counter,number within=section]{methode}[2][]{%
  enhanced,unbreakable,colback=white,colframe=white,boxrule=0pt,arc=0pt,
  borderline north={0.9pt}{0pt}{Bleu},borderline south={0.9pt}{0pt}{Bleu},
  left=4pt,right=4pt,top=8pt,bottom=8pt,before skip=1.4em,after skip=1.4em,
  before upper={\methlabelfix\phantomsection\addcontentsline{met}{metentry}{\protect\makebox[2.6em][l]{\thetcbcounter}#2}%
     {\sffamily\bfseries\small\color{white}\colorbox{Bleu}{\strut\ MÉTHODE~\thetcbcounter\ }}%
     \ \ {\sffamily\bfseries\color{Bleu}#2}\par\smallskip},#1}

% ---------- Repères pédagogiques : un filet de couleur, un mot-clé ----------
\newcommand{\definirbloc}[3]{% nom, couleur, libellé
  \newtcolorbox{#1}[1][]{enhanced,breakable,blanker,
    borderline west={1.4pt}{0pt}{#2},left=10pt,right=2pt,top=1pt,bottom=1pt,
    before skip=1em,after skip=1em,
    before upper={{\sffamily\bfseries\footnotesize\color{#2}\textls[120]{\MakeUppercase{#3}}}\quad},##1}}
\definirbloc{reflexe}{Bleu}{Réflexe}
\definirbloc{astuce}{Ocre}{Astuce}
\definirbloc{attention}{Brique}{Attention}
\definirbloc{enpratique}{BleuMoyen}{En pratique}
\definirbloc{plusloin}{Gris}{Pour aller plus loin}
\definirbloc{strategie}{Bleu}{Stratégie}
\definirbloc{lecture}{BleuMoyen}{Lire le schéma}
% compatibilité avec les parties reprises
\newenvironment{garde}{\begin{attention}}{\end{attention}}
\newenvironment{complementHP}{\begin{plusloin}}{\end{plusloin}}
\newenvironment{interpretation}{\begin{enpratique}}{\end{enpratique}}

% ---------- Exercices, corrigés ----------
\newcounter{exo}[section]
\renewcommand{\theexo}{\thesection.\arabic{exo}}
\newcommand{\niveau}[1]{{\color{Ocre}\footnotesize
  \ifcase#1\or$\bullet\circ\circ\circ\circ$\or$\bullet\bullet\circ\circ\circ$\or$\bullet\bullet\bullet\circ\circ$%
  \or$\bullet\bullet\bullet\bullet\circ$\or$\bullet\bullet\bullet\bullet\bullet$\fi}}
\newcommand{\serie}[1]{\par\needspace{8\baselineskip}\vspace{1.5em}%
  \noindent{\sffamily\bfseries\color{Bleu}#1}\par\nobreak\vspace{-0.4em}%
  \noindent{\color{Filet}\rule{\linewidth}{0.4pt}}\par\nobreak\vspace{0.1em}}
\newenvironment{exo}[3]{%
  \par\needspace{5\baselineskip}\vspace{0.75em}\refstepcounter{exo}\label{e:#1}%
  \def\exolab{#1}%
  \noindent{\sffamily\bfseries\color{Bleu}Exercice~\theexo}\quad\niveau{#2}%
  \hfill{\footnotesize\sffamily\color{Gris}#3}\par\nobreak\vspace{0.15em}}%
  {\par\nopagebreak\noindent\hfill{\scriptsize\sffamily\hyperref[c:\exolab]{\color{Gris}corrigé p.~\pageref{c:\exolab}}}\par}
\newenvironment{corrige}[1]{%
  \par\needspace{6\baselineskip}\vspace{1em}\phantomsection\label{c:#1}%
  \noindent{\sffamily\bfseries\color{Bleu}Exercice~\ref{e:#1}}\quad
  {\scriptsize\sffamily\hyperref[e:#1]{\color{Gris}énoncé p.~\pageref{e:#1}}}\par\nobreak\vspace{0.2em}}%
  {\hfill\textcolor{BleuClair}{$\square$}\par}
\newcommand{\etape}[1]{\par\smallskip\noindent{\sffamily\bfseries\small\color{BleuMoyen}#1.}\ }
\newenvironment{solution}{\par\noindent{\sffamily\bfseries\small\color{BleuMoyen}Solution.}\ }{\hfill\textcolor{BleuClair}{$\square$}\par}
\newcommand{\chercher}{\par\smallskip\noindent{\color{Filet}\rule[0.5ex]{1.2cm}{0.4pt}}\
  {\footnotesize\sffamily\itshape\color{Gris}laisser chercher l'étudiant avant de poursuivre}\par\smallskip}

% ---------- Guide de séance ----------

\newcommand{\marge}[1]{\leavevmode\marginpar{\raggedleft\scriptsize\sffamily\color{Gris}\textls[100]{\MakeUppercase{#1}}}\ignorespaces}
\newcommand{\Etape}[4]{%
  \subsection[Étape #1 --- #3 {\color{Gris}(#2)}]{#3}\label{etape:#1}%
  \noindent{\small\sffamily\color{Gris}{\color{Ocre}\bfseries Étape #1}\quad·\quad #2\quad·\quad #4}\par\medskip}
\newcommand{\question}[1]{\par\smallskip\noindent{\sffamily\bfseries\small\color{BleuMoyen}Question.}\ \textit{#1}\par\smallskip}
\newcommand{\rappel}[1]{\par\smallskip\noindent\phantomsection\addcontentsline{rap}{rapentry}{#1}%
  {\sffamily\bfseries\scriptsize\colorbox{Fond}{\color{Gris}\strut RAPPEL}}\ {\sffamily\small\color{Gris}#1.}\ }
\newcommand{\pourquoi}{\par\noindent{\sffamily\bfseries\small\color{Ocre}Pourquoi ?}\ }

% Ouverture de chapitre : \ouverture{objectifs}{prérequis}{durée}
\newcommand{\ouverture}[3]{%
  \begin{tcolorbox}[enhanced,colback=Fond,colframe=Fond,arc=0pt,boxrule=0pt,left=12pt,right=12pt,top=10pt,bottom=10pt,
     before skip=1.2em,after skip=1.6em]
  \begin{minipage}[t]{0.56\linewidth}
  {\sffamily\bfseries\small\color{Bleu}OBJECTIFS}\par\vspace{0.3em}\small #1
  \end{minipage}\hfill
  \begin{minipage}[t]{0.38\linewidth}
  {\sffamily\bfseries\small\color{Bleu}PRÉREQUIS}\par\vspace{0.3em}\small #2
  \ifx\relax#3\relax\else\par\vspace{0.6em}{\sffamily\bfseries\small\color{Bleu}DURÉE}\par\vspace{0.2em}\small #3\fi
  \end{minipage}
  \end{tcolorbox}}
% Fin de chapitre
\newenvironment{essentiel}{\par\needspace{12\baselineskip}\vspace{1.5em}%
  \begin{tcolorbox}[enhanced,unbreakable,colback=white,colframe=white,boxrule=0pt,arc=0pt,
   borderline north={1.4pt}{0pt}{Ocre},left=2pt,right=2pt,top=8pt,bottom=4pt]
  {\sffamily\bfseries\color{Ocre}L'ESSENTIEL DU CHAPITRE}\par\medskip}{\end{tcolorbox}}

% ---------- Notations ----------
\newcommand{\K}{\mathbb{K}}
\newcommand{\R}{\mathbb{R}}
\newcommand{\C}{\mathbb{C}}
\newcommand{\N}{\mathbb{N}}
\newcommand{\Z}{\mathbb{Z}}
\newcommand{\Mn}[1]{\mathcal{M}_{n}(#1)}
\newcommand{\GLn}[1]{\mathrm{GL}_n(#1)}
\newcommand{\LE}{\mathcal{L}(E)}
\newcommand{\rg}{\operatorname{rg}}
\newcommand{\vect}{\operatorname{Vect}}
\newcommand{\Mat}{\operatorname{Mat}}
\newcommand{\Id}{\operatorname{Id}}
\newcommand{\sh}{\operatorname{sh}}
\newcommand{\ch}{\operatorname{ch}}
\renewcommand{\th}{\operatorname{th}}
\DeclareMathOperator{\argsh}{argsh}
\newcommand{\vv}[1]{\overrightarrow{#1}}
\newcommand{\ve}[1]{\vec{#1}}
\newcommand{\ient}[2]{\llbracket #1,#2\rrbracket}
\newcommand{\dt}[1]{\left|\begin{matrix}#1\end{matrix}\right|}
\newcommand{\dis}{\displaystyle}

% ---------- Figures ----------
\pgfplotsset{trig format plots=rad}
\pgfplotsset{courbe/.style={axis lines=middle,axis equal image,
  axis line style={Gris,-{Stealth[length=5pt]}},
  xlabel={$x$},ylabel={$y$},ticklabel style={font=\scriptsize\sffamily,color=Gris},
  every axis x label/.style={at={(ticklabel* cs:1)},anchor=west,font=\small},
  every axis y label/.style={at={(ticklabel* cs:1)},anchor=south,font=\small},
  samples=300,smooth,width=0.6\linewidth}}
\tikzset{arcetude/.style={Bleu,line width=1.25pt},
         reste/.style={BleuClair,line width=1pt},
         tgte/.style={Ocre,line width=0.9pt},
         asy/.style={Gris,dashed,line width=0.6pt},
         pt/.style={circle,fill=Ink,inner sep=1.2pt}}

% ---------- Navigation ----------
\makeatletter
\newcommand*\l@rapentry{\@dottedtocline{1}{0em}{0em}}
\newcommand*\l@metentry{\@dottedtocline{1}{0em}{0em}}
\newcommand{\listemethodes}{\@starttoc{met}}
\newcommand{\listerappels}{\@starttoc{rap}}
\makeatother
% onglets de chapitre sur le bord droit
\usepackage{eso-pic}
\newif\ifonglets \ongletsfalse
\newcommand{\dessineronglet}{%
  \edef\ongcol{\ongletcouleur}%
  \edef\ongnum{\grandnum}%
  \pgfmathsetmacro{\ongy}{3.2+2.3*(\ongnum-1)}%
  \begin{tikzpicture}[remember picture,overlay]
  \fill[\ongcol] ([xshift=-0.75cm,yshift=-\ongy cm]current page.north east)
     rectangle ([yshift=-\ongy cm-1.9cm]current page.north east);
  \node[text=white,font=\sffamily\bfseries\large] at ([xshift=-0.37cm,yshift=-\ongy cm-0.95cm]current page.north east)
     {\ongnum};
  \end{tikzpicture}}
\AddToShipoutPictureBG{\ifonglets\ifnum\value{section}>0 \dessineronglet\fi\fi}
% blocs de navigation et d'enseignement
\definirbloc{commencer}{Gris}{Si vous commencez ici}
\definirbloc{horsprogramme}{Brique}{Hors programme}
\newcommand{\HP}{\,{\sffamily\bfseries\scriptsize\textcolor{Brique}{[HP]}}}
\newcommand{\source}[1]{{\footnotesize\sffamily\color{Gris}\,(#1)}}
\newcommand{\coupdepouce}[1]{\par\vspace{-0.2em}\noindent{\footnotesize\sffamily\color{Gris}%
  \textbf{Coup de pouce} (à donner si l'étudiant bloque) : #1}\par}
\newcommand{\pointarret}[1]{\par\vspace{1em}\noindent{\color{Filet}\rule[0.5ex]{\linewidth}{0.4pt}}\par\nobreak
  \noindent{\footnotesize\sffamily\color{Ocre}\bfseries POINT D'ARRÊT POSSIBLE}\quad{\footnotesize\sffamily\color{Gris}#1}\par\vspace{0.6em}}
\newcommand{\demarche}[1]{\noindent{\sffamily\bfseries\small\color{BleuMoyen}Démarche.}\ \textit{#1}\par\smallskip}

% ---------- Corrigés : idées, astuces, pièges, à retenir ----------
\newcommand{\choix}[1]{\par\noindent{\sffamily\bfseries\small\color{BleuMoyen}Choix de la méthode.}\ #1\par\smallskip}
\newcommand{\gain}[1]{\par\noindent{\sffamily\bfseries\small\color{Ocre}Astuce.}\ #1\par}
\newcommand{\piege}[1]{\par\noindent{\sffamily\bfseries\small\color{Brique}Piège.}\ #1\par}
\newtcolorbox{retenirbox}{enhanced,breakable,colback=Fond,colframe=Fond,arc=0pt,boxrule=0pt,
  borderline west={2pt}{0pt}{Ocre},left=8pt,right=6pt,top=4pt,bottom=4pt,before skip=0.7em,after skip=0.5em}
\newcommand{\retenir}[1]{\begin{retenirbox}{\sffamily\bfseries\footnotesize\color{Ocre}À RETENIR}\quad #1\end{retenirbox}}
\newcommand{\TD}[1]{{\color{Ocre}\bfseries TD\,#1}\ ·\ }
\newcommand{\q}[1]{\par\smallskip\noindent{\sffamily\bfseries\color{Bleu}#1}\ }
\providecommand{\iddots}{\mathinner{\mkern1mu\raise1pt\hbox{.}\mkern2mu\raise4pt\hbox{.}\mkern2mu\raise7pt\hbox{.}\mkern1mu}}


\providecommand{\Ker}{\operatorname{Ker}}\providecommand{\tr}{\operatorname{tr}}\providecommand{\Sp}{\operatorname{Sp}}
\newcommand{\va}{\vec a}\newcommand{\vb}{\vec b}\newcommand{\vz}{\vec 0}\newcommand{\dd}[1]{\vec\delta_{#1}}
\newcommand{\Qc}{\mathcal Q}
\newcommand{\Etp}[1]{\par\medskip\noindent{\sffamily\bfseries\small\color{Bleu}#1}\enspace\ignorespaces}
\newtcolorbox{etapesbox}{enhanced,breakable,blanker,borderline west={1.6pt}{0pt}{BleuMoyen},left=10pt,right=0pt,top=0pt,bottom=2pt,before skip=0.8em,after skip=0.6em,before upper={\setlength{\parskip}{4pt}}}
\newtcolorbox{enoncebox}{enhanced,breakable,colback=Fond,colframe=Fond,arc=0pt,boxrule=0pt,left=8pt,right=8pt,top=6pt,bottom=6pt,before skip=0.6em,after skip=0.6em}
\newcommand{\Enonce}[1]{\begin{enoncebox}{\sffamily\bfseries\footnotesize\color{Bleu}CE QU'ON DEMANDE}\par\smallskip #1\end{enoncebox}}
\newtcolorbox{ideebox}{enhanced,breakable,colback=BleuClair!18,colframe=BleuClair!18,borderline west={2.4pt}{0pt}{Bleu},arc=0pt,boxrule=0pt,left=10pt,right=8pt,top=6pt,bottom=6pt,before skip=0.6em,after skip=0.8em}
\newcommand{\Idee}[1]{\begin{ideebox}{\sffamily\bfseries\footnotesize\color{Bleu}L'IDÉE}\par\smallskip #1\end{ideebox}}
\newtcolorbox{conclbox}{enhanced,breakable,blanker,borderline west={1.6pt}{0pt}{Ocre},left=8pt,top=2pt,bottom=2pt,before skip=1em,after skip=0.6em}
\newcommand{\Conclusion}[1]{\begin{conclbox}{\sffamily\bfseries\footnotesize\color{Ocre}CONCLUSION}\quad #1\end{conclbox}}
\newtcolorbox{exemplebox}{enhanced,breakable,colback=Ocre!7,colframe=Ocre!7,borderline west={2pt}{0pt}{Ocre},arc=0pt,boxrule=0pt,left=10pt,right=8pt,top=6pt,bottom=6pt,before skip=0.8em,after skip=0.6em}
\newcommand{\Exemple}[1]{\begin{exemplebox}{\sffamily\bfseries\footnotesize\color{Ocre}POUR VOIR, AVEC DES NOMBRES}\par\smallskip #1\end{exemplebox}}
\newcommand{\Explication}{\par\medskip\noindent{\sffamily\bfseries\footnotesize\color{Bleu}L'EXPLICATION}\par}
\newtcolorbox{copie}{enhanced,breakable,colback=white,colframe=BleuMoyen!45,boxrule=0.6pt,arc=2pt,left=10pt,right=10pt,top=8pt,bottom=8pt,before skip=1.1em,after skip=0.8em,
  title={\sffamily\bfseries\footnotesize LA RÉDACTION},attach boxed title to top left={xshift=8pt,yshift=-2.2mm},
  boxed title style={colback=BleuMoyen,colframe=BleuMoyen,arc=1pt,boxrule=0pt,top=1pt,bottom=1pt},coltitle=white,
  before upper={\setlength{\parskip}{5pt}}}
\newcommand{\Redaction}{\par\medskip\noindent{\sffamily\bfseries\footnotesize\color{BleuMoyen}LA RÉDACTION}\par}
\newcommand{\Outil}[2]{\begin{enoncebox}{\sffamily\bfseries\small\color{Bleu}#1}\par\smallskip #2\end{enoncebox}}
\begin{document}
\renewcommand{\grandnum}{1}
\begin{titlepage}
\thispagestyle{empty}
\begin{tikzpicture}[remember picture,overlay]
  \fill[Bleu] (current page.north west) rectangle ([yshift=-15.5cm]current page.north east);
  \node[anchor=north west,text=white,font=\sffamily\small] at ([xshift=2.6cm,yshift=-3.2cm]current page.north west)
     {\textls[250]{MATHÉMATIQUES\ \ ·\ \ CPGE PT\ \ ·\ \ DEUXIÈME ANNÉE}};
  \node[anchor=north west,text=white] at ([xshift=2.5cm,yshift=-4.4cm]current page.north west)
     {\fontsize{62}{66}\selectfont\scshape Le Collé};
  \draw[Ocre,line width=1.6pt] ([xshift=2.6cm,yshift=-7.3cm]current page.north west) -- ++(4.2cm,0);
  \node[anchor=north west,text=white,font=\sffamily\Large\bfseries,align=left]
     at ([xshift=2.6cm,yshift=-8.0cm]current page.north west) {Bases et familles cycliques};
  \node[anchor=north west,text=white!90!Bleu,font=\sffamily\large,align=left]
     at ([xshift=2.6cm,yshift=-9.0cm]current page.north west) {Retour sur le devoir libre 3\\[0.15cm]Endomorphismes cycliques};
  \node[anchor=north west,text=Ink,font=\sffamily\small,align=left,text width=12cm]
     at ([xshift=2.6cm,yshift=-17.4cm]current page.north west)
     {\textcolor{Bleu}{\bfseries Méthodes et exercices corrigés}\\[0.3cm]\textcolor{Gris}{Démontrer qu'une famille est une base, écrire la matrice d'un endomorphisme dans une base bien choisie : les méthodes du DL, pour les prochains devoirs.}};
  \node[anchor=south west,text=Ink,font=\sffamily\large] at ([xshift=2.6cm,yshift=3.0cm]current page.south west) {Ahmed Amine Ouqlani};
  \draw[Filet] ([xshift=2.6cm,yshift=2.6cm]current page.south west) -- ([xshift=-2.4cm,yshift=2.6cm]current page.south east);
\end{tikzpicture}
\end{titlepage}


\thispagestyle{plain}
\section*{Avant-propos}
\markboth{Avant-propos}{}
Ce livret reprend uniquement les questions du DL~3 qui ont posé problème, dans l'ordre du sujet. Presque toutes demandent la
même chose : \emph{démontrer qu'une famille est une base}.

\begin{center}
\small\renewcommand{\arraystretch}{1.3}
\begin{tabular}{@{}l>{\raggedright\arraybackslash}p{8.6cm}r@{}}
\toprule
\textbf{Question} & \textbf{Ce qu'il fallait trouver} & \textbf{Page}\\\midrule
II.4.a & Réutiliser la partie II et traiter les deux sens & \pageref{q:ii4a}\\
II.4.b & Appliquer la partie II à $M^T$ & \pageref{q:ii4b}\\
III.3 & Écrire $g(\va)$, $g^2(\va)$ avec les vecteurs propres & \pageref{q:iii3}\\
III.4 & Calculer $g^3(\va)$ & \pageref{q:iii4}\\
III.5 & Résoudre $(R-\lambda I_3)X=0$ & \pageref{q:iii5}\\
III.6 & Reprendre III.3 avec un vecteur quelconque & \pageref{q:iii6}\\\bottomrule
\end{tabular}
\end{center}

\paragraph{Comment lire} Une question par page, toujours dans le même ordre :
\begin{itemize}[leftmargin=1.2em,itemsep=2pt]
\item \textbf{Ce qu'on demande} : l'énoncé ;
\item \textbf{L'explication} : comment on trouve la preuve, et pourquoi chaque étape ;
\item \textbf{La rédaction} : la réponse telle qu'on l'écrit sur la copie.
\end{itemize}
Les outils de la page~\pageref{q:outils} servent partout. Trois exercices courts, corrigés, terminent le livret.

\vfill
\renewcommand{\contentsname}{Sommaire}
\hypertarget{toc}{}
\tableofcontents
\clearpage
\ongletstrue
% =====================================================================
\section{Bases et familles cycliques}\label{sec:dl}
\renewcommand{\grandnum}{1}

Dans tout le livret, $E$ est de dimension $n$ et $f\in\mathcal L(E)$. On note $f^2=f\circ f$, $f^3=f\circ f\circ f$, etc.

\subsection{Les outils}\label{q:outils}

\Outil{Outil 1 — Base en dimension $n$}{%
Une famille de $n$ vecteurs dans un espace de dimension $n$ est une base dès qu'elle est \emph{libre}. On n'a donc jamais à
prouver qu'elle est génératrice.

Pour prouver qu'elle est libre, on écrit : « soient $c_1,\dots,c_n$ tels que $c_1\vec u_1+\dots+c_n\vec u_n=\vz$ », et on
montre que tous les $c_i$ sont nuls.}

\Outil{Outil 2 — Deux résultats sur la liberté}{%
\begin{itemize}[leftmargin=1.2em,itemsep=2pt]
\item Une famille qui contient $\vz$ est liée : ce n'est jamais une base.
\item Des vecteurs propres associés à des valeurs propres distinctes forment une famille libre.
\end{itemize}}

\Outil{Outil 3 — Endomorphisme ou matrice : c'est pareil}{%
Si $A$ est la matrice de $f$ dans une base, et $X$ la colonne des coordonnées de $\vec x$ dans cette base, alors les
coordonnées de $f(\vec x)$ sont $AX$, et celles de $f^2(\vec x)$ sont $A^2X$.

Dans $\R^3$ avec la base canonique, le vecteur $(x,y,z)$ et sa colonne de coordonnées ont les mêmes nombres. C'est pour cela
qu'on peut calculer $\varphi(\va)$ comme $M\va$, et dire « $\varphi^3=0$ » ou « $M^3=0$ » : c'est la même phrase.
Attention : dans une autre base, la matrice n'est plus $M$.}

\Outil{Outil 4 — « Trouver tous les vecteurs tels que… soit une base »}{%
C'est une équivalence, donc il y a deux choses à prouver : quand la condition est vraie, c'est une base ; quand elle est
fausse, la famille est liée.}

% ---------------------------------------------------------------------
\clearpage
\subsection{II.4.a — Tous les \texorpdfstring{$\va$}{a} qui donnent une base}\label{q:ii4a}

\Enonce{$\varphi$ est l'endomorphisme de $\R^3$ de matrice $M$. Trouver tous les $\va\in\R^3$ tels que
$\Qc=(\va,\varphi(\va),\varphi^2(\va))$ soit une base de $\R^3$.}

\Explication
\begin{etapesbox}
\Etp{Que cherche-t-on ?} On veut savoir pour quels $\va$ la famille $\Qc$ est une base. Ce n'est pas un calcul à faire de
zéro : la question II.1 a déjà démontré que, si $f^3=0$ et $f^2(\va)\neq\vz$, alors $(\va,f(\va),f^2(\va))$ est une base. On va
donc appliquer ce résultat à $\varphi$.

\Etp{Peut-on l'appliquer ?} Il faut d'abord vérifier que $\varphi^3=0$. On le sait grâce à la partie I : $M=PTP^{-1}$, et en
élevant au cube les produits $P^{-1}P$ du milieu se simplifient, donc $M^3=PT^3P^{-1}=0$. Comme $M$ est la matrice de
$\varphi$, cela veut dire $\varphi^3=0$ (outil 3).

\Etp{Pourquoi deux cas ?} « Trouver tous les $\va$ » est une équivalence (outil 4). Si $\varphi^2(\va)\neq\vz$, II.1 conclut. Si
$\varphi^2(\va)=\vz$, le dernier vecteur de $\Qc$ est nul, et une famille qui contient $\vz$ n'est jamais libre (outil 2).

\Etp{Comment rendre la condition concrète ?} On calcule $\varphi^2(\va)=M^2\va$ pour $\va=(x,y,z)$, et on regarde quand ce
vecteur est nul.
\end{etapesbox}

\begin{copie}
D'après la partie I, $M=PTP^{-1}$, donc $M^3=PT^3P^{-1}=0$ : $\varphi^3=0$.

\emph{Si $\varphi^2(\va)\neq\vz$}, d'après II.1, $\Qc$ est une base de $\R^3$.

\emph{Si $\varphi^2(\va)=\vz$}, $\Qc$ contient le vecteur nul, donc elle est liée : ce n'est pas une base.

Pour $\va=(x,y,z)$ :
\[\varphi^2(\va)=M^2\va=\begin{pmatrix}0&0&0\\2&2&-2\\2&2&-2\end{pmatrix}\begin{pmatrix}x\\y\\z\end{pmatrix}
=\begin{pmatrix}0\\2(x+y-z)\\2(x+y-z)\end{pmatrix}.\]
Ainsi $\varphi^2(\va)=\vz\iff x+y-z=0$.

\textbf{Conclusion :} $\Qc$ est une base de $\R^3$ si et seulement si $x+y-z\neq0$.
\end{copie}

% ---------------------------------------------------------------------
\clearpage
\subsection{II.4.b — \texorpdfstring{$M$}{M} est semblable à \texorpdfstring{$M^T$}{tM}}\label{q:ii4b}

\Enonce{Montrer que $M$ est semblable à $M^T$, et que $H=M+2I_3$ est semblable à $H^T$.}

\Explication
\begin{etapesbox}
\Etp{Que veut dire « semblable » ?} $A$ et $B$ sont semblables s'il existe $Q$ inversible avec $B=Q^{-1}AQ$. Il faut donc
\emph{fabriquer} une matrice $Q$ entre $M$ et $M^T$.

\Etp{L'idée en une ligne.} \[M\sim T\quad\text{et}\quad M^T\sim T\qquad\Longrightarrow\qquad M\sim M^T.\]
C'est comme « $a=b$ et $c=b$, donc $a=c$ » : on passe par $T$.

\Etp{Pourquoi $M^T\sim T$ ?} La partie II montre que toute matrice $B$ de taille 3 avec $B^3=0$ et $B^2\neq0$ est semblable
à $T$. On vérifie ces deux conditions pour $M^T$. La transposée d'un produit inverse l'ordre ; pour une puissance, cela ne
change rien : $(M^T)^3=(M^3)^T=0$ et $(M^T)^2=(M^2)^T$. Enfin, la transposée d'une matrice non nulle n'est pas nulle.

\Etp{Comment fabriquer $Q$ ?} On a deux égalités : $M=PTP^{-1}$ et $M^T=STS^{-1}$. De la première, on isole $T$ :
$T=P^{-1}MP$. On le remplace dans la seconde :
\[M^T=S\,\big(P^{-1}MP\big)\,S^{-1}.\]
On pose $Q=PS^{-1}$. L'inverse d'un produit inverse l'ordre : $Q^{-1}=SP^{-1}$. L'égalité s'écrit donc $M^T=Q^{-1}MQ$.

\Etp{Et pour $H$ ?} L'astuce est d'écrire $2I_3=Q^{-1}(2I_3)Q$ : la matrice identité ne change pas quand on la « conjugue ».
La même $Q$ marche donc pour $H=M+2I_3$.
\end{etapesbox}

\begin{copie}
D'après la partie I, $M=PTP^{-1}$, donc $T=P^{-1}MP$.

On a $(M^T)^3=(M^3)^T=0$ et $(M^T)^2=(M^2)^T\neq0$. D'après la partie II, il existe $S$ inversible telle que $M^T=STS^{-1}$.

On remplace $T$ : $M^T=SP^{-1}MPS^{-1}=Q^{-1}MQ$ avec $Q=PS^{-1}$ inversible. Donc $M$ est semblable à $M^T$.

Avec la même matrice $Q$ :
\[H^T=M^T+2I_3=Q^{-1}MQ+2Q^{-1}I_3Q=Q^{-1}(M+2I_3)Q=Q^{-1}HQ.\]
\textbf{Conclusion :} $H$ est semblable à $H^T$.
\end{copie}

% ---------------------------------------------------------------------
\clearpage
\subsection{III.3 — La famille \texorpdfstring{$(\va,g(\va),g^2(\va))$}{(a,g(a),g2(a))} est une base}\label{q:iii3}

\Enonce{$g\in\mathcal L(\C^3)$ a trois valeurs propres distinctes $\lambda_1,\lambda_2,\lambda_3$, de vecteurs propres
$\dd1,\dd2,\dd3$. On pose $\va=\dd1+\dd2+\dd3$. Montrer que $\Qc=(\va,g(\va),g^2(\va))$ est une base de $\C^3$.}

\Explication
\begin{etapesbox}
\Etp{Pourquoi pas comme en II.1 ?} En II.1, on composait par $f^2$ pour tuer des termes, car $f^3=0$. Ici $g$ n'est pas
nilpotent : cette astuce ne marche plus. Il faut une autre idée.

\Etp{Quelle base choisir ?} Dans la base des vecteurs propres $(\dd1,\dd2,\dd3)$, $g$ est très simple : il multiplie chaque
$\dd i$ par $\lambda_i$. On écrit donc $\va$, $g(\va)$, $g^2(\va)$ dans cette base. C'est bien une base : trois vecteurs propres
pour des valeurs propres distinctes sont libres (outil 2), et il y en a 3 en dimension 3 (outil 1).

\Etp{Que faire de la relation ?} On suppose $c_0\va+c_1g(\va)+c_2g^2(\va)=\vz$ et on regroupe tout ce qui multiplie $\dd1$, puis
$\dd2$, puis $\dd3$. Comme les $\dd i$ sont libres, chaque coefficient est nul.

\Etp{D'où vient le polynôme ?} Les trois coefficients ont la même forme $c_0+c_1\lambda+c_2\lambda^2$, avec $\lambda=\lambda_1$,
$\lambda_2$ puis $\lambda_3$. C'est le polynôme $P=c_0+c_1X+c_2X^2$ évalué en trois points. Or un polynôme de degré au plus 2
ne peut pas avoir trois racines distinctes sans être nul.
\end{etapesbox}

\begin{copie}
Les $\dd i$ sont des vecteurs propres associés à des valeurs propres distinctes, donc $(\dd1,\dd2,\dd3)$ est libre ; c'est une
base de $\C^3$. Par linéarité, $g(\va)=\sum_{i=1}^3\lambda_i\dd i$ et $g^2(\va)=\sum_{i=1}^3\lambda_i^2\dd i$.

$\Qc$ a 3 vecteurs et $\dim\C^3=3$ : il suffit de montrer qu'elle est libre. Soient $c_0,c_1,c_2\in\C$ tels que
$c_0\va+c_1g(\va)+c_2g^2(\va)=\vz$. Alors
\[\sum_{i=1}^3\big(c_0+c_1\lambda_i+c_2\lambda_i^2\big)\dd i=\vz.\]
Comme $(\dd i)$ est libre, $P(\lambda_i)=0$ pour $i=1,2,3$, où $P=c_0+c_1X+c_2X^2$. $P$ est de degré au plus 2 et a trois
racines distinctes, donc $P=0$ : $c_0=c_1=c_2=0$.

\textbf{Conclusion :} $\Qc$ est libre, c'est une base de $\C^3$.
\end{copie}

\paragraph{Autre rédaction, plus courte} Dans la base $(\dd i)$, les coordonnées de $\va$, $g(\va)$, $g^2(\va)$ sont
$(1,1,1)$, $(\lambda_1,\lambda_2,\lambda_3)$, $(\lambda_1^2,\lambda_2^2,\lambda_3^2)$. Leur déterminant est un Vandermonde :
\[\det\begin{pmatrix}1&\lambda_1&\lambda_1^2\\1&\lambda_2&\lambda_2^2\\1&\lambda_3&\lambda_3^2\end{pmatrix}
=(\lambda_2-\lambda_1)(\lambda_3-\lambda_1)(\lambda_3-\lambda_2)\neq0.\]

% ---------------------------------------------------------------------
\clearpage
\subsection{III.4 — La matrice de \texorpdfstring{$g$}{g} dans \texorpdfstring{$\Qc$}{Q}}\label{q:iii4}

\Enonce{Montrer que la matrice de $g$ dans $\Qc=(\va,g(\va),g^2(\va))$ est
\[R=\begin{pmatrix}0&0&\alpha\\1&0&\beta\\0&1&\gamma\end{pmatrix},\qquad\text{où }\chi_g(x)=x^3-\gamma x^2-\beta x-\alpha.\]}

\Explication
\begin{etapesbox}
\Etp{Comment lit-on une matrice dans une base ?} La colonne $j$ contient les coordonnées, dans cette base, de l'image du
$j$-ième vecteur de base. Il faut donc calculer $g(\va)$, $g\big(g(\va)\big)$ et $g\big(g^2(\va)\big)$, et les écrire avec
$\va$, $g(\va)$, $g^2(\va)$.

\Etp{Pourquoi les deux premières colonnes sont-elles faciles ?} $g$ envoie chaque vecteur de $\Qc$ sur le suivant :
$g(\va)$ est le 2\textsuperscript{e} vecteur, $g^2(\va)$ le 3\textsuperscript{e}. Leurs coordonnées sont $(0,1,0)$ et $(0,0,1)$.

\Etp{Et la dernière ?} $g^3(\va)$ n'est pas un vecteur de $\Qc$ : il faut l'exprimer. On repasse par les vecteurs propres, comme
en III.3. Chaque $\lambda_i$ est racine de $\chi_g$, ce qui permet de remplacer $\lambda_i^3$ par
$\gamma\lambda_i^2+\beta\lambda_i+\alpha$.
\end{etapesbox}

\begin{copie}
On a $g(\va)=0\cdot\va+1\cdot g(\va)+0\cdot g^2(\va)$ et $g\big(g(\va)\big)=g^2(\va)$ : les deux premières colonnes sont
$(0,1,0)$ et $(0,0,1)$.

Ensuite $g^3(\va)=\sum_{i=1}^3\lambda_i^3\dd i$. Comme $\chi_g(\lambda_i)=0$, on a $\lambda_i^3=\gamma\lambda_i^2+\beta\lambda_i+\alpha$,
donc
\[g^3(\va)=\gamma\sum_{i=1}^3\lambda_i^2\dd i+\beta\sum_{i=1}^3\lambda_i\dd i+\alpha\sum_{i=1}^3\dd i
=\alpha\va+\beta g(\va)+\gamma g^2(\va).\]
\textbf{Conclusion :} la troisième colonne est $(\alpha,\beta,\gamma)$, et la matrice de $g$ dans $\Qc$ est $R$.
\end{copie}

% ---------------------------------------------------------------------
\clearpage
\subsection{III.5 — Les sous-espaces propres dans \texorpdfstring{$\Qc$}{Q}}\label{q:iii5}

\Enonce{Donner les sous-espaces propres de $g$, en coordonnées dans la base $\Qc$.}

\Explication
\begin{etapesbox}
\Etp{Que veut dire « dans la base $\Qc$ » ?} On donne les vecteurs propres par leurs coordonnées dans $\Qc$. Or, dans $\Qc$, la
matrice de $g$ est $R$ (III.4). On résout donc $RX=\lambda_iX$, c'est-à-dire $(R-\lambda_iI_3)X=0$ (outil 3).

\Etp{Par où commencer ?} La dernière ligne du système ne contient que $y$ et $z$ : elle donne $y=(\lambda_i-\gamma)z$. La
deuxième donne alors $x=\lambda_iy-\beta z=(\lambda_i^2-\gamma\lambda_i-\beta)z$.

\Etp{D'où vient la troisième équation ?} Ce n'est pas une nouvelle équation : c'est la première ligne $-\lambda_ix+\alpha z=0$,
dans laquelle on remplace $x$ :
\[-\lambda_i(\lambda_i^2-\gamma\lambda_i-\beta)z+\alpha z=-\big(\lambda_i^3-\gamma\lambda_i^2-\beta\lambda_i-\alpha\big)z=-\chi_g(\lambda_i)\,z.\]
La parenthèse est exactement $\chi_g(\lambda_i)$, puisque $\chi_g(x)=x^3-\gamma x^2-\beta x-\alpha$.

\Etp{Pourquoi est-elle toujours vraie ?} $\lambda_i$ est une valeur propre, donc une racine de $\chi_g$ : $\chi_g(\lambda_i)=0$.
L'équation s'écrit $0\times z=0$, vraie pour tout $z$. Elle n'impose rien : $z$ est libre, et l'espace propre est une droite.
C'était prévisible : si les trois équations imposaient chacune une condition, la seule solution serait $X=0$, alors qu'une valeur
propre a des vecteurs propres non nuls.

\Etp{D'où viennent $\gamma$ et $\beta$ ?} On développe $\chi_g(x)=(x-\lambda_1)(x-\lambda_2)(x-\lambda_3)$ :
\[\chi_g(x)=x^3-(\lambda_1+\lambda_2+\lambda_3)x^2+(\lambda_1\lambda_2+\lambda_1\lambda_3+\lambda_2\lambda_3)x-\lambda_1\lambda_2\lambda_3.\]
On compare avec $x^3-\gamma x^2-\beta x-\alpha$ : $\gamma=\lambda_1+\lambda_2+\lambda_3$ et
$\beta=-(\lambda_1\lambda_2+\lambda_1\lambda_3+\lambda_2\lambda_3)$.

\Etp{Le calcul, pour $i=1$.} On remplace $\gamma$ et $\beta$, et les termes s'annulent deux par deux :
\[\lambda_1-\gamma=\lambda_1-\lambda_1-\lambda_2-\lambda_3=-(\lambda_2+\lambda_3),\]
\[\lambda_1^2-\gamma\lambda_1-\beta=\lambda_1^2-\lambda_1^2-\lambda_1\lambda_2-\lambda_1\lambda_3+\lambda_1\lambda_2+\lambda_1\lambda_3+\lambda_2\lambda_3=\lambda_2\lambda_3.\]
Donc $E_{\lambda_1}=\vect\big((\lambda_2\lambda_3,\,-(\lambda_2+\lambda_3),\,1)\big)$. Pour $\lambda_2$ et $\lambda_3$, c'est le même
calcul en échangeant les rôles : on garde toujours les \emph{deux autres} valeurs propres.
\end{etapesbox}

\Exemple{Avec $\lambda_1=1$, $\lambda_2=2$, $\lambda_3=3$ : $\chi_g(x)=(x-1)(x-2)(x-3)=x^3-6x^2+11x-6$, donc $\gamma=6$,
$\beta=-11$, $\alpha=6$ et
\[R=\begin{pmatrix}0&0&6\\1&0&-11\\0&1&6\end{pmatrix}.\]
Pour $\lambda_1=1$, la formule donne le vecteur $(2\times3,\,-(2+3),\,1)=(6,-5,1)$. On vérifie :
$R\begin{pmatrix}6\\-5\\1\end{pmatrix}=\begin{pmatrix}6\\6+0-11\\-5+6\end{pmatrix}=\begin{pmatrix}6\\-5\\1\end{pmatrix}$ : c'est bien un
vecteur propre pour la valeur propre $1$. Et la première ligne est automatiquement vraie : $-1\times6+6\times1=0$.}

\begin{copie}
D'après III.4, on résout $(R-\lambda_iI_3)X=0$ avec $X=(x,y,z)$ :
\[\begin{cases}-\lambda_ix+\alpha z=0\\ x-\lambda_iy+\beta z=0\\ y+(\gamma-\lambda_i)z=0\end{cases}
\iff\begin{cases}y=(\lambda_i-\gamma)z\\ x=(\lambda_i^2-\gamma\lambda_i-\beta)z\\ -\chi_g(\lambda_i)z=0.\end{cases}\]
La dernière équation est toujours vraie car $\chi_g(\lambda_i)=0$. Donc
$E_{\lambda_i}=\vect\big((\lambda_i^2-\gamma\lambda_i-\beta,\ \lambda_i-\gamma,\ 1)\big)$.

D'après III.1, $\gamma=\lambda_1+\lambda_2+\lambda_3$ et $\beta=-(\lambda_1\lambda_2+\lambda_1\lambda_3+\lambda_2\lambda_3)$. Avec
$\{i,j,k\}=\{1,2,3\}$, on obtient $\lambda_i-\gamma=-(\lambda_j+\lambda_k)$ et $\lambda_i^2-\gamma\lambda_i-\beta=\lambda_j\lambda_k$.

\textbf{Conclusion :} $E_{\lambda_i}=\vect\big((\lambda_j\lambda_k,\ -(\lambda_j+\lambda_k),\ 1)\big)$, en coordonnées dans $\Qc$.
\end{copie}

% ---------------------------------------------------------------------
\clearpage
\subsection{III.6 — Tous les \texorpdfstring{$\vb$}{b} qui donnent une base}\label{q:iii6}

\Enonce{Trouver tous les $\vb\in\C^3$ tels que $(\vb,g(\vb),g^2(\vb))$ soit une base de $\C^3$.}

\Explication
\begin{etapesbox}
\Etp{Quel lien avec III.3 ?} En III.3, $\va=\dd1+\dd2+\dd3$ : toutes ses coordonnées valent $1$. Ici $\vb$ est quelconque :
$\vb=b_1\dd1+b_2\dd2+b_3\dd3$. Comme $g(\dd i)=\lambda_i\dd i$, appliquer $g$ multiplie chaque coordonnée par son $\lambda_i$ :
\[g(\vb)=b_1\lambda_1\dd1+b_2\lambda_2\dd2+b_3\lambda_3\dd3,\qquad g^2(\vb)=b_1\lambda_1^2\dd1+b_2\lambda_2^2\dd2+b_3\lambda_3^2\dd3.\]

\Etp{D'où vient le polynôme ?} On écrit une combinaison et on regroupe selon les $\dd i$ :
\begin{align*}c_0\vb+c_1g(\vb)+c_2g^2(\vb)={}&b_1\big(c_0+c_1\lambda_1+c_2\lambda_1^2\big)\dd1+b_2\big(c_0+c_1\lambda_2+c_2\lambda_2^2\big)\dd2\\&+b_3\big(c_0+c_1\lambda_3+c_2\lambda_3^2\big)\dd3.\end{align*}
La parenthèse a toujours la même forme $c_0+c_1\lambda+c_2\lambda^2$ ; seul $\lambda$ change. On lui donne un nom :
$P=c_0+c_1X+c_2X^2$, et la parenthèse est $P(\lambda_i)$. Le polynôme n'est pas une nouvelle idée : c'est une notation pour
cette parenthèse.

\Etp{Si $b_1$, $b_2$, $b_3$ sont tous non nuls.} On veut prouver que la famille est libre, donc on suppose la combinaison
nulle. Les $\dd i$ sont libres, donc $b_iP(\lambda_i)=0$ pour chaque $i$. On divise par $b_i\neq0$ : $P(\lambda_1)=P(\lambda_2)=
P(\lambda_3)=0$. Un polynôme de degré au plus 2 qui a trois racines est nul : $c_0=c_1=c_2=0$. C'est exactement III.3.

\Etp{Si $b_1=0$ : l'objectif change.} On ne veut plus montrer que c'est libre ; on veut montrer que c'est \emph{lié}. Il faut donc
\emph{trouver} des $c_0,c_1,c_2$ pas tous nuls qui donnent $\vz$. Comme $b_1=0$, le premier terme vaut $0$ quel que soit $P$. Il
suffit donc que $P(\lambda_2)=0$ et $P(\lambda_3)=0$.

\Etp{Comment choisir $P$ ?} On ne résout rien : on fabrique le polynôme qui s'annule là où on veut,
\[P=(X-\lambda_2)(X-\lambda_3)=X^2-(\lambda_2+\lambda_3)X+\lambda_2\lambda_3.\]
On lit ses coefficients : $c_0=\lambda_2\lambda_3$, $c_1=-(\lambda_2+\lambda_3)$, $c_2=1$. Les trois termes sont nuls, et
$c_2=1\neq0$ : c'est une relation non triviale, la famille est liée.

\Etp{L'intuition.} Si $b_1=0$, alors $\vb$, $g(\vb)$, $g^2(\vb)$ n'ont jamais de composante sur $\dd1$ : ils restent tous dans
$\vect(\dd2,\dd3)$, qui est un plan. Trois vecteurs dans un plan sont forcément liés.

\Etp{Pourquoi deux sens ?} C'est encore « trouver tous les vecteurs » : il faut prouver la base dans un cas, et la liaison
dans l'autre (outil 4).
\end{etapesbox}

\Exemple{Avec $\lambda_1=1$, $\lambda_2=2$, $\lambda_3=3$ et $\vb=\dd2+\dd3$ (donc $b_1=0$). En coordonnées dans
$(\dd1,\dd2,\dd3)$ : $\vb=(0,1,1)$, $g(\vb)=(0,2,3)$, $g^2(\vb)=(0,4,9)$. On prend $P=(X-2)(X-3)=X^2-5X+6$, donc
$c_0=6$, $c_1=-5$, $c_2=1$ :
\[6(0,1,1)-5(0,2,3)+(0,4,9)=(0,\ 6-10+4,\ 6-15+9)=(0,0,0).\]
La relation non triviale $6\vb-5g(\vb)+g^2(\vb)=\vz$ montre que la famille est liée.}

\begin{copie}
Soit $\vb=b_1\dd1+b_2\dd2+b_3\dd3$. Par linéarité, $g^k(\vb)=\sum_{i=1}^3b_i\lambda_i^k\dd i$. Pour $c_0,c_1,c_2\in\C$ et
$P=c_0+c_1X+c_2X^2$ :
\[c_0\vb+c_1g(\vb)+c_2g^2(\vb)=\sum_{i=1}^3b_iP(\lambda_i)\dd i.\]

\emph{Si $b_1b_2b_3\neq0$.} Si cette combinaison est nulle, alors $b_iP(\lambda_i)=0$ pour tout $i$ (les $\dd i$ sont libres),
donc $P(\lambda_i)=0$. Comme en III.3, $P=0$ : la famille est libre, c'est une base.

\emph{Si un $b_i$ est nul}, par exemple $b_1=0$. On prend $P=(X-\lambda_2)(X-\lambda_3)$, soit $c_0=\lambda_2\lambda_3$,
$c_1=-(\lambda_2+\lambda_3)$, $c_2=1$. Alors $b_1P(\lambda_1)=0$ et $P(\lambda_2)=P(\lambda_3)=0$, donc la combinaison est
nulle, avec $c_2\neq0$ : la famille est liée.

\textbf{Conclusion :} $(\vb,g(\vb),g^2(\vb))$ est une base si et seulement si $b_1$, $b_2$, $b_3$ sont tous non nuls.
\end{copie}

% ---------------------------------------------------------------------
\clearpage
\subsection{Bilan}\label{q:bilan}

\begin{reflexe}
Pour prouver qu'une famille est libre, regarder la forme des vecteurs :
\begin{center}\small\renewcommand{\arraystretch}{1.3}
\begin{tabular}{@{}>{\raggedright\arraybackslash}p{5.4cm}>{\raggedright\arraybackslash}p{7.2cm}@{}}
\toprule
\textbf{Les vecteurs sont…} & \textbf{On fait…}\\\midrule
des $f^k(\va)$ avec $f^n=0$ & on applique $f^{n-1}$, puis $f^{n-2}$… (II.1)\\
écrits avec des vecteurs propres & on regroupe selon les vecteurs propres, puis polynôme (III.3)\\
des vecteurs numériques & on calcule un déterminant\\\bottomrule
\end{tabular}
\end{center}
\end{reflexe}

\begin{essentiel}
\begin{itemize}[leftmargin=1.1em,itemsep=3pt]
\item $n$ vecteurs en dimension $n$ : il suffit de prouver qu'ils sont libres.
\item « Tous les vecteurs tels que… » : deux sens à prouver.
\item $g^k\big(\sum b_i\dd i\big)=\sum b_i\lambda_i^k\dd i$, et chaque $\lambda_i$ vérifie $\chi_g(\lambda_i)=0$.
\item Endomorphisme et matrice disent la même chose dans une base fixée.
\end{itemize}
\end{essentiel}

\serie{Pour s'entraîner}

\begin{exo}{ii4a}{2}{comme II.4.a}
Soit $f\in\mathcal L(E)$ avec $\dim E=3$, $f^3=0$ et $f^2\neq0$. Montrer que $(\vec x,f(\vec x),f^2(\vec x))$ est une base de $E$ si et
seulement si $\vec x\notin\Ker f^2$.
\end{exo}

\begin{exo}{dim2}{2}{comme III.6}
$(\vec e_1,\vec e_2)$ est une base de $E$ et $f\in\mathcal L(E)$ vérifie $f(\vec e_1)=\vec e_1$, $f(\vec e_2)=-\vec e_2$. Trouver
tous les $\vec u\in E$ tels que $(\vec u,f(\vec u))$ soit une base de $E$.
\end{exo}

\begin{exo}{param}{2}{déterminant}
Pour quelles valeurs de $m\in\R$ les vecteurs $(1,1,m)$, $(1,m,1)$, $(m,1,1)$ forment-ils une base de $\R^3$ ?
\end{exo}

% =====================================================================
\clearpage
\renewcommand{\ongletcouleur}{Ocre}
\section*{Corrigés}\addcontentsline{toc}{section}{Corrigés}\markboth{Corrigés}{}

\begin{corrige}{ii4a}
\emph{Si $\vec x\notin\Ker f^2$.} La famille a $3$ vecteurs en dimension $3$ : on prouve qu'elle est libre. Soit
$c_0\vec x+c_1f(\vec x)+c_2f^2(\vec x)=\vz$. On applique $f^2$ : comme $f^3=f^4=0$, il reste $c_0f^2(\vec x)=\vz$, et
$f^2(\vec x)\neq\vz$, donc $c_0=0$. On applique ensuite $f$ : il reste $c_1f^2(\vec x)=\vz$, donc $c_1=0$. Enfin
$c_2f^2(\vec x)=\vz$ donne $c_2=0$. C'est une base.

\emph{Si $\vec x\in\Ker f^2$}, la famille contient $f^2(\vec x)=\vz$, donc elle est liée.
\end{corrige}

\begin{corrige}{dim2}
On écrit $\vec u=x\vec e_1+y\vec e_2$ ; alors $f(\vec u)=x\vec e_1-y\vec e_2$. En coordonnées,
$\det\begin{pmatrix}x&x\\y&-y\end{pmatrix}=-2xy$. C'est une base si et seulement si $x\neq0$ et $y\neq0$ : exactement comme en
III.6, aucune coordonnée sur les vecteurs propres ne doit être nulle.
\end{corrige}

\begin{corrige}{param}
\[\det\begin{pmatrix}1&1&m\\1&m&1\\m&1&1\end{pmatrix}=-(m-1)^2(m+2).\]
C'est une base si et seulement si $m\neq1$ et $m\neq-2$.
\end{corrige}

\end{document}
